\section*{Hoofdstuk \ref{ch:b1:acetals-protection} --- Carbonylverbindingen: nucleofiele addities, acetalen en bescherming}

\begin{solution}{exo:b1:acetals-protection:1}
\ce{CH3CH(OH)OCH3}: \omterm{def:b1:acetals-protection:hemiacetal}{hemiacetaal}. \ce{CH3CH(OCH3)2}: \omterm{def:b1:acetals-protection:hemiacetal}{acetaal}.
\ce{CH3CH(OH)2}: hydraat. \ce{CH3OCH3}: ether. \ce{(CH3)2C(OCH2CH3)2}:
\omterm{def:b1:acetals-protection:hemiacetal}{acetaal} (van een keton).
\end{solution}

\begin{solution}{exo:b1:acetals-protection:2}
\ce{CH3CH2CHO + 2CH3OH <=> CH3CH2CH(OCH3)2 + H2O}; propanon geeft met
ethaan-1,2-diol 2,2-dimethyl-1,3-dioxolaan en water,
\ce{(CH3)2CO + HOCH2CH2OH <=> C5H10O2 + H2O}.
\end{solution}

\begin{solution}{exo:b1:acetals-protection:3}
Natriumhydroxide, methylmagnesiumbromide en natriumboorhydride laten het
onveranderd; verdund waterig zwavelzuur hydrolyseert het.
\end{solution}

\begin{solution}{exo:b1:acetals-protection:4}
De OH op koolstofatoom 4 addeert aan het aldehydekoolstofatoom 1: een
vijfring van vier koolstofatomen en één zuurstofatoom, met een OH op het
koolstofatoom naast het ringzuurstofatoom.
\end{solution}

\begin{solution}{exo:b1:acetals-protection:5}
Protonering van C=O; methanol addeert aan het koolstofatoom; verlies van
\ce{H+}: \omterm{def:b1:acetals-protection:hemiacetal}{hemiacetaal}; protonering van zijn OH; verlies van water tot
\ce{CH3CH=O+CH3}; een tweede methanol addeert; verlies van \ce{H+}:
\ce{CH3CH(OCH3)2}. De twee opgenomen protonen worden teruggegeven: het
zuur is een \omterm{def:b1:elementary-steps:catalyst}{katalysator}.
\end{solution}

\begin{solution}{exo:b1:acetals-protection:6}
Protonering van één OCH3; verlies van methanol tot \ce{(CH3)2C=O+CH3};
water addeert; verlies van \ce{H+}: \omterm{def:b1:acetals-protection:hemiacetal}{hemiacetaal}; protonering, verlies van
methanol, verlies van \ce{H+}: propanon. Overmaat water houdt $Q$ onder
$K$ voor de hydrolyse, die volledig verloopt.
\end{solution}

\begin{solution}{exo:b1:acetals-protection:7}
Eén water per \omterm{def:b1:acetals-protection:hemiacetal}{acetaal}: $0.250 \times 18.0 = \qty{4.5}{g}$, ongeveer
\qty{4.5}{mL}. Na een uur: $3.2/18.0 = \qty{0.178}{mol}$, een omzetting van
\qty{71}{\percent}.
\end{solution}

\begin{solution}{exo:b1:acetals-protection:8}
Met het diol vervangt één molecuul er twee: de reactie verlaagt het
aantal moleculen niet (één aldehyde en één diol geven één \omterm{def:b1:acetals-protection:hemiacetal}{acetaal} en één
water), en de tweede OH wordt dicht bij het reagerende koolstofatoom
gehouden, wat de ring gemakkelijk sluit. Beide maken het cyclische
\omterm{def:b1:acetals-protection:hemiacetal}{acetaal} gunstiger.
\end{solution}

\begin{solution}{exo:b1:acetals-protection:9}
Het zuur zou het \omterm{def:b1:organomagnesium:organometallic}{grignardreagens} vernietigen (protonoverdracht); ook
sporen water en diol moeten verwijderd worden.
\end{solution}

\begin{solution}{exo:b1:acetals-protection:10}
1. Bescherm het keton van 4-chloorbutaan-2-on als cyclisch \omterm{def:b1:acetals-protection:hemiacetal}{acetaal}
(ethaan-1,2-diol, zure \omterm{def:b1:elementary-steps:catalyst}{katalysator}, Dean--Stark). 2. Magnesium in droge
ether: het \omterm{def:b1:organomagnesium:organometallic}{grignardreagens} van het beschermde chloride, dat niet langer
door zijn eigen carbonylgroep vernietigd wordt. 3. Voeg propanon toe;
opwerking met een waterige oplossing van ammoniumchloride: de tertiaire
alcohol. 4. Verdund waterig zuur: \omterm{def:b1:acetals-protection:protecting-group}{ontscherming} tot
5-hydroxy-5-methylhexaan-2-on.
\end{solution}

\begin{solution}{exo:b1:acetals-protection:11}
$Q = [\text{acetaal}][\ce{H2O}]/([\text{aldehyde}][\ce{CH3OH}]^2)$: een
grote overmaat methanol, gekwadrateerd in de noemer, houdt $Q < K$ tot er
weinig aldehyde overblijft. Water verwijderen werkt op de teller;
methanol als oplosmiddel gebruiken is eenvoudiger, maar vraagt een
product dat gemakkelijk af te scheiden is.
\end{solution}

\begin{solution}{exo:b1:acetals-protection:12}
Alleen de OH van het \omterm{def:b1:acetals-protection:hemiacetal}{hemiacetaal} kan als water vertrekken en een kation
geven dat door het naburige ringzuurstofatoom gestabiliseerd wordt
(\ce{C=O+}); de andere OH-groepen zouden van gewone koolstofatomen moeten
vertrekken, wat niet-gestabiliseerde kationen geeft. Methanol addeert aan
het gestabiliseerde kation: een methylacetaal.
\end{solution}

\begin{solution}{pb:b1:acetals-protection:1}
\textbf{1.} \ce{BrMgCH2CH2CH2CHO}.
\textbf{2.} Zijn carbanionachtige koolstofatoom zou het aldehyde van een
ander molecuul (of van zichzelf) aanvallen: \ce{RMgBr} addeert aan
$\mathrm{R'CHO}$, wat een \omterm{def:b1:alcohol-activation:alkoxide}{alkoxide} geeft.
\textbf{3.} Het aldehyde, tegen het \omterm{def:b1:organomagnesium:organometallic}{grignardreagens}.
\textbf{4.} \omterm{def:b1:acetals-protection:hemiacetal}{Acetalen} zijn stabiel tegenover \omterm{def:b1:organomagnesium:organometallic}{grignardreagentia} en basen en
worden door waterig zuur verwijderd.
\textbf{5.} Het vormt zich vollediger (één diol voor twee
alcoholmoleculen) en is gemakkelijk te hydrolyseren.
\textbf{6.} $15.09/150.9 = \qty{0.1000}{mol}$.
\textbf{7.} $1.20 \times 0.1000 \times 62.0 = \qty{7.44}{g}$.
\textbf{8.} \ce{BrCH2CH2CH2CHO + HOCH2CH2OH <=> C6H11BrO2 + H2O}
(2-(3-broompropyl)-1,3-dioxolaan).
\textbf{9.} Ze activeert de carbonylgroep en maakt het verlies van water
in het tweede stadium mogelijk; ze wordt teruggevormd.
\textbf{10.} Tolueen en water destilleren samen over, condenseren en
vallen in de opvanger; water zakt naar beneden en blijft, tolueen vloeit
terug.
\textbf{11.} $Q$ bevat $[\ce{H2O}]$ in zijn teller; water verwijderen
houdt $Q < K$, en de reactie gaat door.
\textbf{12.} Water is dichter en mengt niet met tolueen: het verzamelt
zich onderaan; alleen de bovenste laag bereikt de terugvoerarm.
\textbf{13.} Wanneer er geen water meer bijkomt, bij het verwachte
volume.
\textbf{14.} $0.1000 \times 194.9 = \qty{19.5}{g}$.
\textbf{15.} \ce{C6H11BrO2 + Mg -> C6H11BrMgO2} (droge ether).
\textbf{16.} Het reagens addeert aan \ce{(CH3)2CO}: het \omterm{def:b1:alcohol-activation:alkoxide}{alkoxide}
\ce{(CH3)2C(O^-)CH2CH2CH2-} op de beschermde keten, met \ce{MgBr+}.
\textbf{17.} Zuur zou de \omterm{def:b1:acetals-protection:protecting-group}{beschermende groep} te vroeg verwijderen en zou
de tertiaire alcohol kunnen dehydrateren.
\textbf{18.} $0.0700 \times 174.0 = \qty{12.2}{g}$.
\textbf{19.} Het \omterm{def:b1:acetals-protection:hemiacetal}{acetaal} + \ce{H2O} $\to$ het aldehyde + ethaan-1,2-diol
(zure \omterm{def:b1:elementary-steps:catalyst}{katalysator}).
\textbf{20.} Haar verlies van water vraagt \omterm{def:b1:predominant-reaction:strong-acid}{sterk zuur} en warmte; mild,
koud waterig zuur hydrolyseert het \omterm{def:b1:acetals-protection:hemiacetal}{acetaal} veel sneller.
\textbf{21.} De OH van koolstofatoom 5 addeert aan het
aldehydekoolstofatoom 1: een zesring (vijf koolstofatomen en één
zuurstofatoom) met twee methylgroepen op het koolstofatoom naast het
ringzuurstofatoom en een OH op het andere koolstofatoom ernaast.
\textbf{22.} $0.0700 \times 0.90 = \qty{0.0630}{mol}$, \qty{8.19}{g}.
\textbf{23.} $0.0630/0.1000 = \qty{63}{\percent}$ (\omterm{def:b1:acetals-protection:protecting-group}{bescherming} als
volledig beschouwd).
\textbf{24.} $0.1000 \times 18.0 = \qty{1.80}{g}$.
\textbf{25.} $1.80/0.997 =$ \textbf{\qty{1.8}{mL} water} in de opvanger.
\end{solution}
